% =============================================================================
% Semana 08 — Trabalho 4 + Storytelling: Big Idea e Storyboard
% Aula de 3 horas (19h–22h) | 20/04/2026 — somente segunda-feira
% =============================================================================

\chapter{Semana 08 (20/04/2026) — Storytelling com Dados: Big Idea e Storyboard}

% ---------------------------------------------------------------------------
\section{Objetivo da semana}
% ---------------------------------------------------------------------------

Pessoal, bem-vindos à Unidade II! A Unidade I nos deu as ferramentas para
\textbf{entender} os dados. Agora a pergunta central muda:

\begin{center}
\textit{``Como eu conto essa história para que alguém tome uma decisão?''}
\end{center}

Análise exploratória serve ao analista. Análise explanatória serve ao
\textbf{tomador de decisão}. Hoje aprendemos a diferença e os primeiros
blocos de qualquer apresentação de dados eficaz: a \textbf{Big Idea} e
o \textbf{Storyboard}.

\begin{NoteBox}
\textbf{Calendário:} aula concentrada — somente segunda-feira (20/04).
Roteiro adaptado para 3 horas incluindo apresentações do T4.
\end{NoteBox}

\begin{BoardBox}
\textbf{Roteiro da aula (19h–22h)}
\begin{enumerate}
  \item 19h–20h — \textbf{Apresentações:} Trabalho 4 (primeiro contato com Tableau)
  \item 20h–20h20 — \textbf{TEORIA:} AED vs Explanatória — a virada de perspectiva
  \item 20h20–21h — \textbf{TEORIA:} Big Idea Framework (Nancy Duarte)
  \item 21h–21h30 — \textbf{TEORIA:} Estrutura narrativa e Storyboard
  \item 21h30–21h50 — \textbf{PRÁTICA:} Escrever Big Idea + montar storyboard de 3 telas
  \item 21h50–22h — \textbf{Correção coletiva + lançamento do T5}
\end{enumerate}
\end{BoardBox}

\begin{NoteBox}
\textbf{Livro de referência para a Unidade II:} \textit{Datastory} de Nancy Duarte
(2019). Capítulos 1–3 para esta semana.
\end{NoteBox}

% =============================================================================
\section{Parte I — Apresentações: Trabalho 4 (19h–20h)}
% =============================================================================

\begin{BoardBox}
\textbf{Checklist de avaliação T4 — Primeiro contato com Tableau:}
\begin{tabular}{|p{4.5cm}|p{8.5cm}|}
\hline
\textbf{Ponto} & \textbf{O que observar} \\
\hline
Dataset importado & Dados carregados corretamente no Tableau \\
Pelo menos 1 visual & Gráfico funcional e com título significativo \\
Título que é mensagem & Não ``Gráfico de Barras'' — sim ``Produto X lidera vendas em SP'' \\
Big Idea presente & 1 frase com ponto de vista + contexto \\
Wireframe/esboço & 3 telas planejadas (mesmo manual/papel) \\
\hline
\end{tabular}
\end{BoardBox}

% =============================================================================
\section{Parte II — Teoria: AED vs Explanatória (20h–20h20)}
% =============================================================================

\begin{SolvedBox}
\textbf{A virada de perspectiva:}

\textbf{Análise Exploratória (AED):}
\begin{itemize}
  \item Você como analista explorando os dados.
  \item Objetivo: entender o dataset, encontrar padrões, detectar anomalias.
  \item Centenas de visuais gerados; a maioria descartada.
  \item Linguagem: ``O dado mostra $X$.''
  \item Audiência: você mesmo.
\end{itemize}

\textbf{Análise Explanatória:}
\begin{itemize}
  \item Você como comunicador levando uma descoberta para o tomador de decisão.
  \item Objetivo: mudar uma crença ou provocar uma ação específica.
  \item Apenas os visuais que servem à \emph{mensagem} são apresentados.
  \item Linguagem: ``Como este dado muda o que você deveria fazer?''
  \item Audiência: CEO, gerente, cliente, investidor.
\end{itemize}

\textbf{Erro mais comum de analistas:} mostrar todo o processo exploratório
para uma audiência executiva. O resultado: confusão, perda de tempo, e a
sua descoberta mais importante se perde no meio de 30 gráficos.
\end{SolvedBox}

% =============================================================================
\section{Parte III — Big Idea Framework (20h20–21h)}
% =============================================================================

\begin{BoardBox}
\textbf{O que é a Big Idea? (para o quadro)}

A Big Idea é uma única frase que captura a essência da sua mensagem.
Ela deve conter \textbf{três elementos obrigatórios}:

\begin{enumerate}
  \item \textbf{Ponto de vista:} uma afirmação, não uma observação neutra.
  (``O dado mostra X'' não é Big Idea. ``Devemos fazer Y porque X'' é
  Big Idea.)
  \item \textbf{Tensão:} o que está em jogo? Qual o risco de não agir?
  Qual a oportunidade desperdiçada?
  \item \textbf{Stake:} por que a audiência deveria se importar? O que a
  afeta diretamente?
\end{enumerate}
\end{BoardBox}

\begin{SolvedBox}
\textbf{Big Ideas boas vs. ruins:}

\begin{tabular}{|p{6cm}|p{7cm}|}
\hline
\textbf{Ruim (observação sem direção)} & \textbf{Boa (Big Idea com tensão)} \\
\hline
``A taxa de cancelamento subiu 12\% no Q3.'' &
``Se não corrigirmos o onboarding até setembro, perderemos R\$1,2M em receita
recorrente nos próximos 12 meses.'' \\
\hline
``Clientes Premium têm NPS mais alto.'' &
``Convertendo 15\% dos clientes Standard para Premium podemos aumentar o NPS
médio em 8 pontos e reduzir o churn em 22\%.'' \\
\hline
``O canal online representa 60\% das vendas.'' &
``O canal físico gera 40\% das vendas com 70\% do custo — precisamos decidir
sua continuidade antes do próximo trimestre.'' \\
\hline
\end{tabular}

\textbf{Teste da Big Idea:} leia para alguém que não viu os dados.
Se a resposta for ``interessante'' sem que a pessoa saiba o que fazer,
a Big Idea não está pronta.
\end{SolvedBox}

\begin{BoardBox}
\textbf{A Pergunta de Decisão (complemento da Big Idea)}

Toda boa apresentação de dados responde a uma \textbf{pergunta de decisão}:
uma questão que a audiência precisa responder \emph{ao final da apresentação}.

\textbf{Exemplos de Perguntas de Decisão:}
\begin{itemize}
  \item ``Devemos expandir para o canal digital ou reforçar o físico?''
  \item ``Qual produto descontinuamos no próximo trimestre?''
  \item ``Aprovamos o orçamento para o novo sistema de CRM?''
\end{itemize}

O analista não precisa ter a resposta — mas precisa apresentar os dados que
tornam a resposta possível.
\end{BoardBox}

% =============================================================================
\section{Parte IV — Estrutura Narrativa e Storyboard (21h–21h30)}
% =============================================================================

\begin{SolvedBox}
\textbf{Estrutura narrativa clássica adaptada para dados (Nancy Duarte):}

\begin{enumerate}
  \item \textbf{Contexto — O que é:} mostre o estado atual. Dados reais.
  A situação é X. A audiência reconhece esse mundo.
  \item \textbf{Conflito — O que poderia ser:} introduza a tensão. O mundo
  \emph{poderia} ser Y — mas algo impede isso. Ou: o mundo \emph{está indo} para Z,
  e isso é um problema.
  \item \textbf{Resolução — O que proponho:} sua proposta de ação.
  Com base nos dados, recomendo W.
  \item \textbf{Call to action — O que peço:} qual decisão específica
  você quer da audiência até quando?
\end{enumerate}

\textbf{Em uma apresentação de dados, isso se traduz em:}
\begin{itemize}
  \item Tela 1: contexto (KPIs atuais, cenário)
  \item Tela 2: problema/oportunidade (o dado que revela a tensão)
  \item Tela 3: análise/evidências
  \item Tela 4: recomendação + próximo passo concreto
\end{itemize}
\end{SolvedBox}

\begin{BoardBox}
\textbf{Storyboard: planejando antes de abrir o Tableau (para o quadro)}

Storyboard é o rascunho da apresentação — desenhado no papel ou em qualquer
ferramenta simples — antes de qualquer visual ser produzido.

\textbf{Por que fazer storyboard antes?}
\begin{itemize}
  \item Força a pensar na \emph{mensagem} antes de pensar no gráfico.
  \item Evita criar visuais que não servem à narrativa.
  \item Alinha expectativas com a audiência sem custo de produção.
\end{itemize}

\textbf{Cada tela do storyboard deve ter:}
\begin{enumerate}
  \item Título que é mensagem (ex.: ``Canal físico custa 4× mais por venda'')
  \item Tipo de visual proposto (barras, linha, mapa, etc.)
  \item O que o dado vai mostrar (esboço manual)
  \item A transição lógica (``por isso, na próxima tela, veja...'')
\end{enumerate}
\end{BoardBox}

% =============================================================================
\section{Parte V — Prática: Big Idea + Storyboard (21h30–21h50)}
% =============================================================================

\begin{SolvedBox}
\textbf{Cenário:} você é analista de dados de uma rede de academias.
O dataset mostra:
\begin{itemize}
  \item Taxa de retenção de alunos: \textbf{62\%}
  \item Taxa de retenção de alunos que usam o app: \textbf{89\%}
  \item Apenas \textbf{28\%} dos alunos instalaram o app
  \item Custo de aquisição de novo aluno: R\$ 180
  \item Custo de retenção via campanha de app: R\$ 12 por aluno
\end{itemize}

\textbf{Tarefa:}
\begin{enumerate}
  \item Escreva a Big Idea (1 frase: ponto de vista + tensão + stake).
  \item Escreva a Pergunta de Decisão.
  \item Monte um storyboard de 3–4 telas (pode ser manual):
  título de cada tela + tipo de visual + transição.
\end{enumerate}
\textbf{Tempo:} 20 minutos. Individual.
\end{SolvedBox}

\begin{SolvedBox}
\textbf{Gabarito orientativo (professor):}

\textbf{Big Idea:} ``Aumentar a adoção do app de 28\% para 70\% pode elevar
a taxa de retenção de 62\% para mais de 80\% — com custo 15× menor do que
adquirir novos alunos para repor os que cancelam.''

\textbf{Pergunta de Decisão:} ``Aprovamos a campanha de adoção do app no próximo mês?''

\textbf{Storyboard:}
\begin{enumerate}
  \item \textbf{Tela 1} — ``62\% de retenção: estamos perdendo 38 em cada 100 alunos ao ano''
  (visual: KPI grande + linha temporal). Transição: ``Mas há um grupo diferente...''
  \item \textbf{Tela 2} — ``Quem usa o app fica: 89\% de retenção vs 62\%''
  (visual: barras comparativas, destacando a diferença de 27pp). Transição: ``O problema: poucos usam.''
  \item \textbf{Tela 3} — ``Apenas 28\% instalaram o app'' (visual: gauge ou pizza impactante).
  Transição: ``E o custo de mudar isso é irrisório.''
  \item \textbf{Tela 4} — ``Campanha de app: R\$12/aluno vs R\$180 para substituir um que cancela''
  (visual: comparação de barras custo ×). \textbf{Call to action:} aprovação até 30/04.
\end{enumerate}
\end{SolvedBox}

% ---------------------------------------------------------------------------
\section{Lançamento — Trabalho 5}
% ---------------------------------------------------------------------------

\begin{SolvedBox}
\textbf{Trabalho 5 — Storytelling com Dados (entrega: 29/04/2026 às 23h59):}
\begin{itemize}
  \item[$\square$] Big Idea (1 frase completa com ponto de vista, tensão e stake)
  \item[$\square$] Pergunta de Decisão explícita
  \item[$\square$] Storyboard de 4–6 telas (manual ou digital)
  \item[$\square$] Pelo menos 2 visuais no Tableau
  \item[$\square$] Títulos que são mensagens (não rótulos de tipo de gráfico)
  \item[$\square$] Script de apresentação de 3 min (por escrito)
\end{itemize}
\end{SolvedBox}

% ---------------------------------------------------------------------------
\section{Materiais Preparados}
% ---------------------------------------------------------------------------
\begin{itemize}
  \item Slides: comparativo AED vs Explanatória; fórmula visual da Big Idea
  \item Folha de storyboard em branco (4 retângulos por folha — impressa para cada aluno)
  \item Exemplos de Big Ideas boas e ruins (slides com quiz rápido)
  \item Transcrição do TED Talk de Nancy Duarte ``The secret structure of great talks''
\end{itemize}

% ---------------------------------------------------------------------------
\section{Próximo passo}
% ---------------------------------------------------------------------------

Na \textbf{Semana 09} aprofundamos o \textbf{design de informação}: o Teste do
Relance, princípios de pré-atenção (cor, tamanho, posição), título que carrega
a mensagem e como transformar visuais ruins em comunicação clara.

Perguntem aos alunos: \textit{``Olhando para o storyboard que vocês fizeram
hoje, qual das 4 telas é a mais importante? Se vocês tivessem apenas 30 segundos
para apresentar — qual tela mostrariam e por quê?''}
